\documentclass{beamer}
\usepackage{minted} % For code listing
\usepackage{fancyvrb}

% Useful macros
% Display C/C++ codes
\newenvironment{cppcode}[2]
%[language=c++,caption=#2,numbers=left,numbersep=5pt,listingname=Code,label=#1,style=colorful]
{\VerbatimEnvironment \begin{minted}{C++}}
{\end{minted}}

\newcommand{\cppinline}[1]



\usetheme{CambridgeUS}
\usefonttheme{structurebold}

% Presentation metadata
\title[Use template]{Using this template to create slides}
\subtitle[Using beamer]{Using beamer}
\author{Balajiganapathi S}
\institute{code-drills.com}
\logo{\includegraphics[height=0.25cm]{code-drills.png}}
\date{\today} % Auto substitutes date. Freeze after completion

% Before each section show outline with curent section bolded and only its subsections
\AtBeginSection[]{\begin{frame}\frametitle{Outline}\tableofcontents[currentsection,hideothersubsections]\end{frame}}

% Begin slides
\begin{document}

% Title page displays metadata
\frame{\titlepage{}}

% Table of contents of whole presentation; shows only section title
\begin{frame}
    \frametitle{Outline}
    \tableofcontents[hideallsubsections]
\end{frame}

% Start a section/subsection
\section{Getting started}
\subsection{Setting up structure}

% First frame: demonstrates pause and newlines
\begin{frame}
    \frametitle{Copy template directory}
    Copy paste the template directory \pause{} \\
    Then modify the contents as you please \pause{} \\
    Use pdflatex to parse and output the corresponding pdf \\
\end{frame}

% List items demo
\subsection{Steps to do after setting up}
\begin{frame}
    \frametitle{Basic details}
    \begin{itemize}
        \item Change name of template file to proper lecture name
        \item Fill in the metadata $-$ title, subtitle etc.
        \item Modify the frames
    \end{itemize}
\end{frame}

% Multiple slides in a subsection
\begin{frame}
    \frametitle{Creating slides}
    See existing sections and subsections for examples and modify it.
\end{frame}

% Demos alert block, graphics
\section{License}
\subsection{Conditions for use}

\begin{frame}
    \frametitle{Creative Commons}
    All \mintinline{C++}|bfs(competitive_programming)| slides are licensed under \includegraphics[height=0.5cm]{CC-BY-SA_icon.png}
    \begin{alertblock}{CC-BY-SA Restrictions}
        \alert{BY} $-$ Licensees may copy, distribute, display and perform the work and make derivative works and remixes based on it only if they give the author or licensor the credits (attribution) in the manner specified by these. \\
        \alert{SA} $-$ Licensees may distribute derivative works only under a license identical (''not more restrictive'') to the license that governs the original work.  \\
    \end{alertblock}
\end{frame}

\section{More demos}
\begin{frame}[fragile]
    \frametitle{Code listing}

    \begin{cppcode}{helloworld}{Hello World}
        #include<iostream>
        
        using namespace std;

        int main() {
            string name = "Balajiganapathi S";
            int rank = 1;
            double score = 100.0;
            cout << "Hello, World!" << endl;
            return 0;
        }
    \end{cppcode}
\end{frame}

% Conclusion section, include references
\section{Conclusion}
\begin{frame}
    \frametitle{Thank you}
    Well begin is half done
\end{frame}

\end{document}
